% Reusable rule-view TikZ module
% Usage: \RuleModule{<visibility>}{<view robots: color/x/y list>}{<center color>}{<direction>}{<dir label>}{<view label>}
% Directions: left, right, top, btm, idle
% Example: \RuleModule{1}{redbot/1/0}{bluebot}{right}{R}{Rule 1}
\newcommand{\RuleModule}[6]{%
  \begin{tikzpicture}[scale=0.68,>=Stealth,line cap=round]
    \def\V{#1}
    \begin{scope}
      \clip (0,\V) -- (\V,0) -- (0,-\V) -- (-\V,0) -- cycle;
      \foreach \x in {-\V,...,\V} {%
        \draw[densely dashed,gray!70] (\x,-\V) -- (\x,\V);
      }
      \foreach \y in {-\V,...,\V} {%
        \draw[densely dashed,gray!70] (-\V,\y) -- (\V,\y);
      }
    \end{scope}

    \foreach \x in {-\V,...,\V} {%
      \foreach \y in {-\V,...,\V} {%
        \pgfmathtruncatemacro{\d}{abs(\x) + abs(\y)}
        \ifnum\d>\V\relax\else
          \fill[gray!25] (\x,\y) circle (0.12);
        \fi
      }
    }

    % Robots visible in the current view.
    \foreach \c/\rx/\ry in {#2} {%
      \pgfmathparse{\rx}\let\rrx\pgfmathresult
      \pgfmathparse{\ry}\let\rry\pgfmathresult
      \filldraw[\c,draw=black] (\rrx,\rry) circle (0.18);
    }

    % Central robot that applies the rule.
    \filldraw[#3,draw=black] (0,0) circle (0.18);

    % Direction action is applied to the central robot only.
    \def\Dir{#4}
    \def\Left{left}
    \def\Right{right}
    \def\Top{top}
    \def\Btm{btm}
    \def\Idle{idle}
    \ifx\Dir\Left
      \draw[->,thick,#3] (-0.24,0) -- (-0.6,0)
        node[midway,above,#3,font=\scriptsize] {#5};
    \fi
    \ifx\Dir\Right
      \draw[->,thick,#3] (0.24,0) -- (0.6,0)
        node[midway,above,#3,font=\scriptsize] {#5};
    \fi
    \ifx\Dir\Top
      \draw[->,thick,#3] (0,0.24) -- (0,0.6)
        node[midway,left,#3,font=\scriptsize] {#5};
    \fi
    \ifx\Dir\Btm
      \draw[->,thick,#3] (0,-0.24) -- (0,-0.6)
        node[midway,left,#3,font=\scriptsize] {#5};
    \fi
    \ifx\Dir\Idle
      \draw[->,thick,#3] (50:0.18) to[out=60,in=0,looseness=6] (-50:0.18)
        node[midway,above right=10pt,#3,font=\scriptsize] {#5};
    \fi
    % Optional bottom-center view label
    \def\ViewLabel{#6}
    \ifx\ViewLabel\empty\else
      \node[font=\scriptsize,black!70,anchor=north] at (0,-\V-0.18) {#6};
    \fi
  \end{tikzpicture}%
}

% View-only module (no direction, no color indicator)
% Usage: \ViewModule{<visibility>}{<view robots: color/x/y list>}{<center color>}{<label>}
% Example: \ViewModule{1}{redbot/1/0}{bluebot}{}
\newcommand{\ViewModule}[4]{%
  \begin{tikzpicture}[scale=0.68,>=Stealth,line cap=round]
    \def\V{#1}
    \begin{scope}
      \clip (0,\V) -- (\V,0) -- (0,-\V) -- (-\V,0) -- cycle;
      \foreach \x in {-\V,...,\V} {%
        \draw[densely dashed,gray!70] (\x,-\V) -- (\x,\V);
      }
      \foreach \y in {-\V,...,\V} {%
        \draw[densely dashed,gray!70] (-\V,\y) -- (\V,\y);
      }
    \end{scope}
    \foreach \x in {-\V,...,\V} {%
      \foreach \y in {-\V,...,\V} {%
        \pgfmathtruncatemacro{\d}{abs(\x) + abs(\y)}
        \ifnum\d>\V\relax\else
          \fill[gray!25] (\x,\y) circle (0.12);
        \fi
      }
    }
    \foreach \c/\rx/\ry in {#2} {%
      \pgfmathparse{\rx}\let\rrx\pgfmathresult
      \pgfmathparse{\ry}\let\rry\pgfmathresult
      \filldraw[\c,draw=black] (\rrx,\rry) circle (0.18);
    }
    \filldraw[#3,draw=black] (0,0) circle (0.18);
    \def\ViewLabel{#4}
    \ifx\ViewLabel\empty\else
      \node[font=\scriptsize,black!70,anchor=north] at (0,-\V-0.18) {#4};
    \fi
  \end{tikzpicture}%
}

% Parallel rule module (two active moving robots with their own directions and colors)
% Usage: \ParallelRuleModule{<visibility>}{<bot1 color>}{<bot1 dir>}{<bot1 label>}{<bot2 x>}{<bot2 y>}{<bot2 color>}{<bot2 dir>}{<bot2 label>}
\newcommand{\ParallelRuleModule}[9]{%
  \begin{tikzpicture}[scale=0.68,>=Stealth,line cap=round]
    \def\V{#1}
    
    % Dashed grid segments (only inside each robot's visibility cross)
    % Horizontal segments
    \draw[densely dashed,gray!70] (-1,0) -- (2,0);
    
    % Vertical segments
    \draw[densely dashed,gray!70] (0,-1) -- (0,1);
    \draw[densely dashed,gray!70] (1,-1) -- (1,1);

    % Gray background circles ONLY at the 8 visible cells (Manhattan distance <= 1 from either center)
    \fill[gray!25] (-1,0) circle (0.12);
    \fill[gray!25] (0,0) circle (0.12);
    \fill[gray!25] (1,0) circle (0.12);
    \fill[gray!25] (2,0) circle (0.12);
    \fill[gray!25] (0,1) circle (0.12);
    \fill[gray!25] (1,1) circle (0.12);
    \fill[gray!25] (0,-1) circle (0.12);
    \fill[gray!25] (1,-1) circle (0.12);

    % Central robot 1 (at 0,0)
    \filldraw[#2,draw=black] (0,0) circle (0.18);

    % Obstacle (at -1,0)
    \filldraw[colorO,draw=black] (-1,0) circle (0.18);

    % Central robot 2 (at #5,#6)
    \filldraw[#7,draw=black] (#5,#6) circle (0.18);

    % Bot 1 Movement arrow
    \def\DirOne{#3}
    \def\Left{left}
    \def\Right{right}
    \def\Top{top}
    \def\Btm{btm}
    \ifx\DirOne\Left
      \draw[->,thick,#2] (-0.24,0) -- (-0.6,0) node[midway,above,#2,font=\scriptsize] {#4};
    \fi
    \ifx\DirOne\Right
      \draw[->,thick,#2] (0.24,0) -- (0.6,0) node[midway,above,#2,font=\scriptsize] {#4};
    \fi
    \ifx\DirOne\Top
      \draw[->,thick,#2] (0,0.24) -- (0,0.6) node[midway,left,#2,font=\scriptsize] {#4};
    \fi
    \ifx\DirOne\Btm
      \draw[->,thick,#2] (0,-0.24) -- (0,-0.6) node[midway,left,#2,font=\scriptsize] {#4};
    \fi

    % Bot 2 Movement arrow (relative to #5,#6)
    \def\DirTwo{#8}
    \ifx\DirTwo\Left
      \draw[->,thick,#7] (#5-0.24,#6) -- (#5-0.6,#6) node[midway,above,#7,font=\scriptsize] {#9};
    \fi
    \ifx\DirTwo\Right
      \draw[->,thick,#7] (#5+0.24,#6) -- (#5+0.6,#6) node[midway,above,#7,font=\scriptsize] {#9};
    \fi
    \ifx\DirTwo\Top
      \draw[->,thick,#7] (#5,#6+0.24) -- (#5,#6+0.6) node[midway,left,#7,font=\scriptsize] {#9};
    \fi
    \ifx\DirTwo\Btm
      \draw[->,thick,#7] (#5,#6-0.24) -- (#5,#6-0.6) node[midway,left,#7,font=\scriptsize] {#9};
    \fi

    % Bottom-center label
    \node[font=\scriptsize,black!70,anchor=north] at (0.5,-1.3) {P-Rule};
  \end{tikzpicture}%
}
